%% World-Model Calculus for Probabilistic Logic Networks
%%
%% Render with:
%%   tectonic world-model-calculus.tex
%% Or:
%%   pdflatex world-model-calculus.tex && bibtex world-model-calculus \
%%   && pdflatex world-model-calculus.tex && pdflatex world-model-calculus.tex

\documentclass[]{article}
\usepackage[utf8]{inputenc}
\usepackage{url}
\usepackage[hyperindex,breaklinks]{hyperref}
%\usepackage{breakurl}  % incompatible with pdflatex/tectonic
\usepackage{listings}
\lstset{
  basicstyle=\ttfamily\footnotesize,
  breaklines=true,
  frame=single,
  extendedchars=false,
  literate={->}{$\to$}{2} {<->}{$\leftrightarrow$}{3}
           {forall}{$\forall$}{6} {exists}{$\exists$}{6}
}
\usepackage{float}
\restylefloat{table}
\usepackage{longtable}
\usepackage{graphicx}
\usepackage[font=small,labelfont=bf]{caption}
\usepackage{amsfonts}
\usepackage{amssymb}
\usepackage{amsmath}
\usepackage{amsthm}
\usepackage{mathtools}
\usepackage{enumitem}
\usepackage{booktabs}

\newtheorem{theorem}{Theorem}
\newtheorem{lemma}[theorem]{Lemma}
\newtheorem{definition}[theorem]{Definition}
\newtheorem{remark}[theorem]{Remark}
\newtheorem{example}[theorem]{Example}
\newtheorem{corollary}[theorem]{Corollary}

\newcommand{\code}[1]{\texttt{#1}}
\newcommand{\Evidence}{\mathsf{Evidence}}
\newcommand{\State}{\mathsf{State}}
\newcommand{\Query}{\mathsf{Query}}
\newcommand{\WM}{\mathsf{WM}}
\newcommand{\revise}{\mathbin{\oplus}}
\newcommand{\tensor}{\mathbin{\otimes}}
\newcommand{\Strength}{\mathsf{strength}}
\newcommand{\ev}{\mathsf{evidence}}
\newcommand{\side}{\Sigma}

\title{World-Model Calculus for Probabilistic Logic Networks:\\
  Evidence, Views, and $\Sigma$-Guarded Compiled Inference\\
  \texttt{(DRAFT --- \today)}}
\author{Zar\footnote{Corresponding author.}
\and Oru\v{z}i\footnote{``Oru\v{z}i'' denotes AI-assisted collaboration
using Claude Code (Anthropic, Opus 4.6), Codex (OpenAI, GPT-5.*), and
ChatGPT Pro (GPT-5.*). Co-authored under human direction.}}

\begin{document}
\maketitle

\begin{abstract}
We present a \emph{World-Model (WM) calculus} for Probabilistic Logic
Networks (PLN), formalized in Lean~4 as part of the \code{Mettapedia}
repository.

The guiding principle is: \emph{PLN inference is stateful posterior revision
plus evidence extraction}, and ``fast PLN rules'' are \emph{compiled,
$\Sigma$-guarded rewrites} whose soundness is discharged relative to a
declared \emph{world-model class} (Bayesian network, factor graph, joint
evidence, etc.).

We separate:
(i)~\textbf{evidence} as a rich algebraic object (quantale/Heyting-valued),
(ii)~\textbf{views} (probabilities, STVs/WTVs, typicalities) as explicit
projections from evidence, and
(iii)~\textbf{tractable sublayers} where inference is exact relative to
assumptions (e.g.\ d-separation).
We explain how this architecture supports message passing / activation
spreading, how to make approximation risk explicit, and how to integrate
quantifiers via satisfying-set semantics and Goertzel-style weakness.
We show that the WM calculus is not merely an abstract evidence API: it
can faithfully host both Kripke and neighborhood modal semantics as
concrete instantiations, recovering ordinary pointed-model truth on
singleton states and yielding an evidential ensemble semantics on
multisets. The neighborhood instance admits the strongest bridge: an
implication-level soundness/completeness theorem connecting proof-theoretic
derivability to WM consequence inequalities. A partial governance-fragment
translation into neighborhood semantics is also provided.

Finally, we outline soundness and (fragmented) completeness results
relative to chosen semantics.

\paragraph{Status note.}
Lean-proven claims are those tied to explicit Lean file/theorem references
(primarily \S\ref{sec:artifacts}). Other sections are research synthesis,
design rationale, or roadmap. Verify against current build outputs before
citing as formalization state.
\end{abstract}

\tableofcontents
\newpage

%% ============================================================
\section{Introduction}
\label{sec:intro}
%% ============================================================

PLN aims to support \emph{probabilistic reasoning in open-ended domains}:
uncertain facts, imperfect independence, and large-scale inference control.
The historical friction is well-known:

\begin{itemize}[leftmargin=2em]
\item \textbf{Local truth-value rules are not globally complete.}
Even with local evidence for $A,B,C$ and links $A\Rightarrow B$,
$B\Rightarrow C$, the \emph{exact} evidence for $A\Rightarrow C$ cannot in
general be computed without joint correlation information.
\item \textbf{Independence assumptions are the real price.}
When is a Markov/d-separation assumption justified? What happens when it is
violated? How do you quantify drift~\cite{Pearl1988}?
\item \textbf{Operational factor graphs are not automatically semantics.}
A message-passing graph whose factors are ``deduction'' or ``induction''
update formulas is useful for search, but semantically correct only when
those factors are \emph{compiled} from a true world-model factorization (or
guarded by explicit assumptions).
\end{itemize}

The design answer:

\begin{quote}
\textbf{Keep the semantic kernel small and honest.}
World-model states revise under evidence, and queries extract evidence.
Everything else---truth values, probabilities, ``fast rules'', message
passing---are \emph{explicit views and $\Sigma$-guarded tactics} whose
validity is proved \emph{relative to a declared model class}.
\end{quote}

This aligns several foundational perspectives:
Kolmogorov/Cox/Knuth--Skilling give valuation axioms and regraduation
principles~\cite{KnuthSkilling2012},
Kohlas--Shenoy give algebraic local computation~\cite{KohlasShenoy2000},
and Meredith/Stay want the categorical/process structure to stay
compositional.


%% ============================================================
\section{The WM Kernel: States, Revision, and Queries}
\label{sec:kernel}
%% ============================================================

\subsection{World-model interface}

At the semantic core we posit:
\begin{itemize}[leftmargin=2em]
\item a type of world-model states $W:\State$,
\item a type of queries $q:\Query$,
\item an evidence type $\Evidence$, and
\item operations \emph{revision} and \emph{evidence extraction}.
\end{itemize}

Conceptually:
\[
  W \revise \Delta \;\;\text{is revision of state $W$ by new evidence
  $\Delta$,} \qquad
  \ev(W,q)\in \Evidence \;\;\text{is evidence for $q$ in state $W$.}
\]

In Lean this is packaged as a \code{WorldModel} typeclass with an
\code{EvidenceType} for states. The essential law is \textbf{additivity of
evidence extraction under revision}:
\[
  \ev(W_1 \revise W_2, q) = \ev(W_1,q) \revise \ev(W_2,q).
\]
This is the ``semantic spine'' that makes revision and querying behave like
a calculus.

For the concrete \code{PathMap}/\code{Finset}$\,\alpha$ instance, this law
is Lean-proven as \code{pjoin\_evidence\_additive} in
\code{OSLF/PathMap/PLNBridge.lean}: for disjoint stores $W_1 \perp W_2$,
\[
  \ev(W_1 \cup W_2, q) = \ev(W_1,q) \oplus \ev(W_2,q),
\]
where $\ev(W,q) = (|W \cap q|,\, |W \setminus q|)$. The proof also
establishes the K\&S sum rule ($n^+ + n^- = |W|$), monotonicity, the
Bayesian conditioning partition, and the STV connection
($\Strength = |W\cap q|/|W|$). Axioms: standard Lean kernel only.

A \code{Multiset}$\,\alpha$ variant (\code{WorldModelBridge.lean})
resolves the idempotency obstacle: \code{Finset} union is idempotent
($W \cup W = W$), breaking \code{evidence\_add} for non-disjoint pairs.
\code{Multiset} sum is the free commutative monoid, enabling proper
non-idempotent revision.


\subsection{Judgments as a proof-calculus skeleton}
\label{sec:judgments}

The WM kernel yields sequent-style judgments:

\begin{itemize}[leftmargin=2em]
\item $\vdash_{\WM} W$ --- \emph{$W$ is a derivable WM state},
\item $\vdash_q\, W \Downarrow q \mapsto e$ --- \emph{from state $W$,
  query $q$ yields evidence $e$},
\item $\vdash_s\, W \Downarrow q \mapsto s$ --- a scalar \emph{strength
  view}.
\end{itemize}

This is already a \emph{proof calculus}, but a \emph{stateful} one: proof
objects are states and evidence terms, not only formulas.

Three judgment variants track provenance at increasing precision:

\begin{enumerate}[leftmargin=2em]
\item \code{WMJudgment W} --- context-free: ``$W$ is derivable.''
\item \code{WMJudgmentCtx}~$\Gamma$~$W$ --- set-indexed: ``$W$ is
  derivable from states in $\Gamma$.'' Tightens the trust boundary.
\item \code{WMJudgmentMulti}~$\Gamma$~$W$ --- multiset-indexed: tracks
  multiplicities. Records if the same observation was used twice.
\end{enumerate}


\subsection{Query equivalence and $\Sigma$-guarded rewrites}
\label{sec:rewrites}

\begin{definition}[Query equivalence]
Two queries $q_1,q_2$ are \emph{evidence-equivalent} if
\[
  \forall W.\;\; \ev(W,q_1)=\ev(W,q_2).
\]
\end{definition}

A ``PLN inference rule'' is packaged as a sound rewrite:

\begin{definition}[$\Sigma$-guarded rewrite rule]
\label{def:rewrite}
A rewrite rule consists of:
\begin{itemize}[leftmargin=2em]
\item a side condition $\side$ (a proposition),
\item a conclusion query $q$,
\item a derived evidence term $\mathrm{derive}(W)$, and
\item a proof that $\side \Rightarrow \forall W.\;
  \mathrm{derive}(W)=\ev(W,q)$.
\end{itemize}
\end{definition}

Many PLN rules rewrite only the \emph{strength view}: they prove an
identity about $\Strength(\ev(W,q))$ rather than about full evidence.

Query equivalence admits three levels of refinement, all formalized:
\begin{itemize}[leftmargin=2em]
\item \code{WMQueryEq} --- full evidence equality (strongest).
\item \code{WMEvidenceLE} --- pointwise dominance (preorder).
  Antisymmetry recovers \code{WMQueryEq}.
\item \code{WMViewEq} --- agreement after projecting through a view
  function $v : \Evidence \to X$ (weakest).
\end{itemize}


%% ============================================================
\section{Evidence vs Views}
\label{sec:evidence}
%% ============================================================

\subsection{Evidence as an algebraic object}

PLN needs an evidence object that encodes positive and negative support,
partial truth / paraconsistency, monotone aggregation, and (in some
sublayers) conjugate-prior counts.

The concrete type is:
\[
  \Evidence = (n^+, n^-) \in \mathbb{R}_{\ge 0}^\infty \times
  \mathbb{R}_{\ge 0}^\infty.
\]
This forms:
\begin{itemize}[leftmargin=2em]
\item a \textbf{complete lattice} (coordinatewise $\le$),
\item a \textbf{commutative monoid} under $\oplus$ (coordinatewise
  addition: evidence aggregation),
\item a \textbf{commutative monoid} under $\otimes$ (coordinatewise
  product: sequential composition),
\item a \textbf{quantale}: $\otimes$ distributes over arbitrary joins.
\end{itemize}

All algebraic laws are Lean-proven in \code{EvidenceQuantale.lean}.


\subsection{Views as projections}

A \textbf{view} is a function $v:\Evidence\to X$, chosen for operational
convenience:
\[
  \text{strength}\quad s(e) \in [0,1], \qquad
  \text{confidence}\quad c(e) \in [0,1], \qquad
  \text{ITV}\quad [l,u] \subseteq [0,1].
\]
Views \emph{must be explicit}, because they necessarily forget information.
This enforcement prevents silent unsoundness.

\begin{remark}[The totality gate]
If $\Evidence$ has incomparable elements, there is no faithful
order-embedding into $\mathbb{R}$. So a single real-valued view cannot be
globally complete. This is not a bug; it is a statement about what
information you are discarding.
\end{remark}

The view layer is parameterized by \code{InterpretableEvidence}:
\begin{itemize}[leftmargin=2em]
\item \code{BinaryContext} (Beta priors) --- Bayesian strength/confidence.
\item \code{ContinuousContext} (Normal-Gamma) --- continuous inference.
\item \code{CategoricalContext} (Dirichlet) --- multinomial inference.
\end{itemize}

ITV (Indefinite Truth Value) semantics provide credible intervals on both
Bayesian (Beta posterior) and Walley IDM
branches~\cite{Walley1996BagOfMarbles}.


%% ============================================================
\section{No-Go Theorem}
\label{sec:nogo}
%% ============================================================

A foundational result proved in Lean:

\begin{theorem}[No local complete link-deduction rule]
\label{thm:nogo}
There is no function
\[
  f:\Evidence^5\to \Evidence
\]
that takes only local link/proposition evidence (for $A,B,C$ and
$A\Rightarrow B$, $B\Rightarrow C$) and returns the exact evidence for
$A\Rightarrow C$ for all world-model states.
\end{theorem}

\noindent
Interpretation: \textbf{correlations live in the joint state.} ``Fast PLN
rules'' must be framed as admissible rewrites under assumptions; they are
not globally complete. This no-go is the reason the WM calculus is
\emph{state-first}: the state carries whatever correlation structure is
needed.

See \code{PLNJointEvidenceNoGo.lean} for the Lean proof.


%% ============================================================
\section{Tractable Sublayers}
\label{sec:sublayers}
%% ============================================================

\subsection{World-model classes and $\Sigma$ side conditions}

To scale, we restrict the class of admissible states:
\[
  \mathcal{C} \subseteq \State
  \quad \text{(e.g.\ DAG-factorizing BNs, Markov random fields, etc.)}.
\]
Inside such a class, additional theorems (screening-off, d-separation) hold
and can be compiled into admissible rewrite rules~\cite{Pearl1988}.

Side conditions $\side$ include:
\begin{itemize}[leftmargin=2em]
\item structural hypotheses (a particular DAG/factorization),
\item conditional independence (from d-separation),
\item positivity/nondegeneracy (to avoid division by $0$),
\item approximation/error ledger invariants.
\end{itemize}


\subsection{Variable elimination as exact query engine}

Given a discrete factor graph (or BN compiled to a factor graph), we
compute unnormalized weights by variable elimination (VE)~\cite{Dechter1999}:
\[
  Z(\text{constraints}) \;=\; \sum_{\text{consistent assignments}}
  \prod_f \phi_f.
\]
In Lean, a generic VE engine is implemented for factor graphs parameterized
by a semiring $K$, with specialization to $\mathbb{R}_{\ge 0}^\infty$.

\begin{remark}
A \code{CommSemiring} assumption is \emph{stronger} than separate
\code{AddCommMonoid}+\code{Mul}, but it is not ``making it a ring.''
Rings add additive inverses; semirings do not. VE is ``sum of products''
and needs distributivity.
\end{remark}


\subsection{Fast PLN rules as compiled tactics}

A PLN rule (e.g.\ deduction) becomes:

\begin{quote}
A $\Sigma$-guarded rewrite of a hard query into smaller queries, together
with a proof that the strength view of the conclusion equals the rule
formula applied to the strength views of the premises.
\end{quote}

This is a \emph{strength-level} rewrite unless the model class supplies a
law of evidence flow. The evidence for the conclusion is still obtained by
querying the WM state, unless stronger results hold.

Concrete instances include BN screening-off
(\code{PLNBNCompilation.lean}) and fast Bayesian-network
rules (\code{PLNBayesNetFastRules.lean}).


%% ============================================================
\section{Knuth--Skilling Bridges}
\label{sec:ks}
%% ============================================================

Knuth--Skilling (K\&S) ``foundations of inference'' axiomatize valuation
schemes on lattices~\cite{KnuthSkilling2012}. In our setting, K\&S plays
two roles:

\begin{enumerate}[leftmargin=2em]
\item \textbf{Probability as a valuation/view.}
Within Boolean world sets, probabilities arise as valuations.
\item \textbf{Regraduation bridges for computation.}
K\&S representation theorems justify mapping an abstract valuation to a
real scale where combination becomes addition:
\[
  \Theta(x \star y) = \Theta(x) + \Theta(y).
\]
This is explicitly a \textbf{scale change}, not probability normalization.
\end{enumerate}

This matters for PLN because it provides a principled bridge between
abstract evidence/valuation objects and numerical algorithms; it lets the
VE engine operate over K\&S-valued factor graphs after regrading; and it
keeps ``probability'' as a \emph{view} rather than the foundational type.

The Lean formalization lives in
\code{KnuthSkilling/Bridges/ValuationAlgebra.lean}.


%% ============================================================
\section{Factor Graphs: Semantics vs Control}
\label{sec:fg}
%% ============================================================

A key conceptual cleanup:

\begin{enumerate}[leftmargin=2em]
\item \textbf{Semantic factor graphs:} factors \emph{are the model}.
They represent a genuine factorization (BN/MRF), and VE / junction tree is
exact relative to that model class.
\item \textbf{Operational factor graphs:} factors are rule schemas or
control edges. Message passing is a heuristic that schedules which queries
to ask. It becomes \emph{semantics} only after compilation: each
operational factor is justified by a semantic factor (or by a proven
admissible rewrite under $\Sigma$).
\end{enumerate}


\subsection{Evidence ledgers and provenance}

Message passing, activation spreading, and positive feedback loops are not
disallowed. They need accountability:

\begin{itemize}[leftmargin=2em]
\item Track provenance and overlap to avoid double-counting evidence.
\item Track $\Sigma$ hypotheses explicitly.
\item Maintain a ledger of approximation error / drift estimates.
\end{itemize}

Evidence sources can overlap (shared observations, ``data incest'').
A principled fix is a \textbf{partial} commutative monoid (separation
algebra): $W_1 \revise W_2$ is defined only when fragments are known
independent/disjoint. If disjointness is uncertain, either keep fragments
separate (delayed fusion) or fuse conservatively using
bounds~\cite{Walley1996BagOfMarbles}.


%% ============================================================
\section{Quantifiers in the WM Calculus}
\label{sec:quantifiers}
%% ============================================================

PLN must support probabilistic reasoning over quantified claims.

\subsection{Satisfying sets and weakness}

A clean PLN-compatible semantics treats a predicate as a frame-valued
satisfying set:
\[
  \chi_P : U \to \Evidence,
\]
analogous to a subobject classifier in a
topos~\cite{MacLaneMoerdijk1992}. Universal quantification is evaluated
via Goertzel's \emph{weakness} on the diagonal relation:
\[
  \forall x.\,P(x)
  \quad\leadsto\quad
  \mathrm{weakness}_\mu\big(\{(u,v)\mid P(u)\wedge P(v)\}\big).
\]

Two distinct quantifier views are kept separate:
\begin{itemize}[leftmargin=2em]
\item \textbf{Extensional (classical) quantifiers:} ``all individuals satisfy.''
\item \textbf{Typicality quantifiers:} ``generality'' measured by
  weakness/diagonal mass.
\end{itemize}

The separation prevents subtle category mistakes: typicality is not
classical $\forall$, but it is exactly what many PLN uses want.

In the WM architecture, quantified statements are simply queries $q$: the
world model state decides how to compute $\ev(W,q)$. Different
world-model classes yield different quantifier semantics: finite-domain
exact, sampled approximate, or quasi-Borel/infinite.

See \code{PLNFirstOrder/QuantifierSemantics.lean}.

\subsection{Algorithmic Quantifier Theorem Pack (Ch.~11)}

\code{PLNFirstOrder/QuantifierAlgorithmTheorems.lean} (361 lines, 0~sorries)
proves three families of results about the algorithmic fuzzy-quantifier
layer:

\paragraph{Parameter monotonicity.}
The near-one, near-zero, and near-one-fraction scores are monotone in
$\varepsilon$ (\code{nearOne\_mono\_of\_epsilon\_le},
\code{nearZero\_mono\_of\_epsilon\_le},
\code{nearOneFraction\_mono\_of\_epsilon\_le}).  Relaxing $\varepsilon$
or~PCL weakens the for-all threshold
(\code{fuzzyForAllHolds\_of\_epsilon\_and\_PCL\_relax},
\code{fuzzyIntervalHolds\_of\_wider\_bounds}).

\paragraph{Crisp endpoints.}
When $\mathrm{PCL}=1$ the fuzzy for-all collapses to a fraction-equals-one
test (\code{fuzzyForAllHolds\_iff\_fraction\_eq\_one\_of\_PCL\_eq\_one}),
and the fuzzy there-exists collapses to a near-zero-fraction-equals-zero
test (\code{fuzzyThereExistsHolds\_iff\_nearZeroFraction\_eq\_zero\_of\_PCL\_eq\_one}).
At the fully crisp endpoint $\varepsilon=0$ the for-all is equivalent to
all scores being~1 (\code{crispForAll\_endpoint\_iff\_allEqOne}).

\paragraph{Rule-4 discrepancy bounds.}
Rule-4 measures the gap between conjunction-based and quantification-based
scoring.  In the \emph{high-gate} regime
($g\ge$ all individual scores), the discrepancy vanishes when the gate
covers all witnesses
(\code{rule4Discrepancy\_bound\_of\_gate\_high},
\code{rule4Discrepancy\_eq\_zero\_of\_gate\_high\_and\_cover}).
In the \emph{low-gate} regime ($g<$ all scores), the discrepancy
equals the minimum score and is bounded by~$g$
(\code{rule4Discrepancy\_eq\_minScore\_of\_gate\_low},
\code{rule4Discrepancy\_le\_g\_of\_gate\_low}).


\subsection{Higher-order quantification}

Quantifying over predicates is hard in crisp logics.
Probabilistic semantics changes the game:

\begin{itemize}[leftmargin=2em]
\item \textbf{HO probabilities can be flattened.}
Kyburg argues that second-order probabilities add no conceptual
power~\cite{Kyburg1988}: they can be represented as ordinary joint
distributions over expanded spaces.
\item \textbf{Henkin semantics as a pragmatic HO target.}
Quantify only over a chosen collection of predicates, yielding
completeness metatheorems analogous to FO settings.
\item \textbf{Quasi-Borel / measurable HO as the long game.}
HO quantification can be interpreted as integration over function
spaces---difficult, but mathematically well-defined.
\end{itemize}

\subsection{Third-Order Quantifier Semantics}

\code{PLNFirstOrder/ThirdOrderQuantifierSemantics.lean} (200 lines,
0~sorries) makes the ``distribution over second-order uncertainty'' layer
explicit.  Each domain element carries a \code{SecondOrderUncertainty}
(a probability distribution over latent outcomes, each with an
\code{Evidence} value); the collection over the domain forms a
\code{ThirdOrderQuantifierModel}.

The key results:
\begin{itemize}[leftmargin=2em]
\item Universal evidence from the full domain bounds observed evidence
  (\code{forAllEvidenceUniversal\_le\_observed}), with equality when all
  elements are observed
  (\code{forAllEvidenceObserved\_univ\_eq\_universal}).
\item Universal confidence is bounded by finite-observation confidence
  (\code{confidence\_universal\_le\_observed}).
\item A \emph{weakness bridge} connects the third-order model to the
  existing \code{forAllEval} semantics under explicit pointwise
  assumptions
  (\code{forAllEvalWeakness\_eq\_forAllEval\_of\_assumptions},
  \code{forAllEvalWeakness\_confidence\_eq\_of\_assumptions}).
\end{itemize}


%% ============================================================
\section{Connections}
\label{sec:connections}
%% ============================================================

\subsection{Temporal/fixpoint reasoning}

Once reasoning becomes stateful (WM revision) and iterative (message
passing), temporal/fixpoint logics become natural. The modal
$\mu$-calculus provides a bridge: temporal operators (``eventually'',
``always'') can be presented as least/greatest fixpoints; PLN temporal
operators embed under a suitable semantics; and fixpoint structure aligns
with recurrent ``anchor states'' and invariant tracking.

\subsection{Process calculi and linear logic}

WM states can be treated as \emph{resources} (linear logic flavor).
Revision is a commutative monoidal addition of evidence/resources. Query
answering is a form of observation (a morphism out of state). Operational
factor graphs become process networks, while the WM semantics provides
correctness conditions.

\subsection{Intensional world models}

The intensional layer (\code{PLNIntensionalWorldModel.lean}) extends
queries with sort-indexed channels: extensional, intensional-ASSOC,
intensional-PAT, and mixed inheritance. This is packaged as
\code{WorldModelSigma} --- a typed variant where queries form a dependent
family $\Query : \mathsf{Srt} \to \mathsf{Type}$.


%% ============================================================
\section{What ``Sound and Complete'' Should Mean Here}
\label{sec:metatheory}
%% ============================================================

\subsection{Soundness}

For a fragment $\mathcal{L}$ and semantics $[\![\cdot]\!]$:
\begin{itemize}[leftmargin=2em]
\item \textbf{Kernel soundness:} revision/query laws (additivity).
\item \textbf{Rewrite soundness:} each $\Sigma$-guarded rule preserves
  the chosen view.
\item \textbf{Approximate soundness:} rules preserve bounds or carry
  explicit error budget.
\end{itemize}

\subsection{Completeness (fragmented, but meaningful)}

Completeness is stated relative to a declared semantics and fragment:
\begin{itemize}[leftmargin=2em]
\item \textbf{Propositional:} evidence-quantale semantics vs
  IPL/intermediate logics.
\item \textbf{First-order:} satisfying-set/weakness quantifiers (finite
  or Henkin FO).
\item \textbf{BN fragment:} completeness relative to BN query answering
  (VE/junction tree).
\item \textbf{Higher-order:} Henkin completeness; QBS as aspirational.
\end{itemize}

This is not a retreat. It is how serious proof theory works: different
fragments have different theorems.

\paragraph{Scope of HOL/FOL completeness wrappers.}
The infinitary HOL and FOL completeness modules
(\code{PLNWorldModelHOLInfinitaryCompleteness.lean},
\code{PLNWorldModelFOLInfinitaryCompleteness.lean}) provide
\emph{implication-level closure wrappers} over declared
soundness/completeness assumptions of the underlying proof systems: given
a pointwise implication $\varphi \Rightarrow \psi$ (or its singleton-strength
equivalent), they produce a \code{WMConsequenceRule} witnessing
$\Strength(W,\varphi) \le \Strength(W,\psi)$ for all multiset states~$W$.
These are not global WM-shell completeness results; they transport
proof-system guarantees into WM strength inequalities.


%% ============================================================
\section{Lean Artifact Map}
\label{sec:artifacts}
%% ============================================================

The main Lean artifacts (all in \code{Mettapedia/}):

\begin{longtable}{p{0.55\textwidth}p{0.4\textwidth}}
\hline
\textbf{File} & \textbf{Role} \\ \hline
\endhead
\code{Logic/EvidenceClass.lean} & Evidence typeclasses,
  \code{InterpretableEvidence}, contexts \\
\code{Logic/EvidenceQuantale.lean} & \code{Evidence} structure, complete
  lattice, quantale, views \\
\code{Logic/PLNWorldModel.lean} & WM kernel: \code{WorldModel},
  \code{PLNQuery}, \code{WMJudgment} \\
\code{Logic/PLNWorldModelCalculus.lean} & $\Sigma$-guarded rewrites,
  \code{WMQueryEq}, strength rules \\
\code{Logic/PLNJointEvidence.lean} & \code{JointEvidence} (Dirichlet over
  worlds) \\
\code{Logic/PLNJointEvidenceNoGo.lean} & No-go theorem
  (Theorem~\ref{thm:nogo}) \\
\code{Logic/PLNBayesNetWorldModel.lean} & BN sublayer: CPT evidence,
  marginalization \\
\code{Logic/PLNBNCompilation.lean} & Screening-off rewrites for BN \\
\code{Logic/PLNBayesNetFastRules.lean} & Compiled fast BN rules \\
\code{Logic/PLNIntensionalWorldModel.lean} & Sort-indexed queries,
  \code{WorldModelSigma} \\
\code{Logic/PLNWorldModelInstitution.lean} & Institution-style signature
  transport with explicit satisfaction-condition theorems \\
\code{Logic/PLNWorldModelHyperdoctrine.lean} & Hyperdoctrine-style typed
  query fibers, change-of-base, Beck-Chevalley transport \\
\code{Logic/PLNWorldModelCategoricalBridge.lean} & Unified endpoint
  theorem linking institution transport and Beck-Chevalley transport \\
\code{Logic/PLNWorldModelITV.lean} & ITV semantics (Bayes and
  Walley~IDM) \\
\code{Logic/PLNWorldModelITVHypercube.lean} & Interval-selector layer,
  typed ITV transport \\
\code{Logic/PLNCanonicalAPI.lean} & One-call end-to-end endpoints \\
\code{Logic/PLNFirstOrder/QuantifierSemantics.lean} & Extensional vs
  typicality quantifiers \\
\code{Logic/PLNFirstOrder/QuantifierAlgorithmTheorems.lean} &
  Ch.~11 algorithmic quantifier monotonicity, crisp endpoints, rule-4
  discrepancy bounds (361 lines) \\
\code{Logic/PLNFirstOrder/ThirdOrderQuantifierSemantics.lean} &
  Third-order distribution model, weakness bridge,
  confidence-degradation bound (200 lines) \\
\code{ProbabilityTheory/BayesianNetworks/FactorGraph.lean} & Factor graph
  semantics \\
\code{ProbabilityTheory/BayesianNetworks/VariableElimination.lean} & VE
  query engine \\
\code{ProbabilityTheory/KnuthSkilling/Bridges/ValuationAlgebra.lean} &
  K\&S regraduation bridge \\
\code{OSLF/PathMap/PLNBridge.lean} & PathMap $\leftrightarrow$ PLN
  evidence bridge \\
\code{OSLF/PathMap/WorldModelBridge.lean} & Multiset WM instance \\
\code{Logic/LP/WorldModelBridge.lean} & LP least Herbrand model bridge \\
\code{OSLF/Framework/QuantaleCoherence.lean} & Quantale
  transport/coherence \\
\code{Logic/PLNWorldModelKripke.lean} & Kripke WM instance,
  singleton adequacy \\
\code{Logic/PLNWorldModelKripkeConsequence.lean} & Deontic Kripke
  consequence rules \\
\code{Logic/PLNWorldModelKripkeCompleteness.lean} & Kripke
  sound/complete proof-theoretic bridge \\
\code{Logic/PLNWorldModelNeighborhood.lean} & Neighborhood WM instance,
  singleton adequacy \\
\code{Logic/PLNWorldModelNeighborhoodCompleteness.lean} &
  Sound/complete proof-theoretic bridge \\
\code{Logic/PLNWorldModelNeighborhoodConsequence.lean} & Governance
  translation, modal families \\
\code{Logic/PLNWorldModelKripkeNeighborhoodEmbedding.lean} &
  Generic Kripke-to-neighborhood embedding \\
\code{Logic/PLNWorldModelKripkeNeighborhoodCanonical.lean} &
  Canonical representation bridge \\
\code{Logic/PLNWorldModelKripkeWeighted.lean} &
  Weighted/trusted-gate multiset semantics \\
\code{Logic/PLNWorldModelExperiment.lean} &
  Deterministic experiment/channel WM interface + Blackwell consequence wrapper \\
\code{Logic/PLNWorldModelExperimentRegression.lean} &
  Finite deterministic experiment fixture regression \\
\code{Logic/PLNWorldModelExperimentStochastic.lean} &
  Stochastic-channel Blackwell utility layer + weighted/trusted-source transfer \\
\code{Logic/PLNWorldModelExperimentStochasticRegression.lean} &
  Finite stochastic Blackwell/weighted-policy fixture regression \\
\code{Logic/PLNWorldModelHOLInfinitary.lean} &
  Countable $\bigwedge$/$\bigvee$ HOL WM instance, singleton adequacy,
  pointwise-implication strength transfer (291 lines) \\
\code{Logic/PLNWorldModelHOLInfinitaryCompleteness.lean} &
  Implication-closure wrappers: \code{WMConsequenceRule} from
  pointwise/singleton-strength (89 lines) \\
\code{Logic/PLNWorldModelFOLInfinitary.lean} &
  Countable $\bigwedge$/$\bigvee$ FOL WM instance, singleton adequacy,
  pointwise-implication strength transfer (297 lines) \\
\code{Logic/PLNWorldModelFOLInfinitaryCompleteness.lean} &
  Implication-closure wrappers: \code{WMConsequenceRule} from
  pointwise/singleton-strength (93 lines) \\
\code{Logic/PLNWorldModelInfinitaryRegression.lean} &
  Concrete HOL+FOL infinitary fixture endpoints (129 lines) \\
\hline
\end{longtable}

All files build with zero sorries and standard Lean kernel axioms only
(\code{propext}, \code{Classical.choice}, \code{Quot.sound}).


%% ============================================================
\section{Modal Semantics Instantiations}
\label{sec:modal}
%% ============================================================

The generic WM kernel (\S\ref{sec:kernel}) is a semantic evaluation
framework: every state is derivable, and query judgments deterministically
extract evidence. The question is whether this framework can host
\emph{real} modal semantics---not merely as illustration, but with genuine
adequacy and proof-theoretic closure. This section describes nine Lean
files (2{,}913 lines total, zero sorries) that answer affirmatively.


\subsection{Kripke WM instance}
\label{sec:kripke-wm}

A \code{PointedKripke} structure pairs a Kripke model (frame + valuation)
with a designated world. A multiset of such structures forms a WM state,
and evidence is extracted by counting satisfaction and refutation across
the ensemble:
\[
  \ev(\{m_1,\ldots,m_n\},\,\varphi)
  = \bigl(|\{i \mid m_i \models \varphi\}|,\;
         |\{i \mid m_i \not\models \varphi\}|\bigr).
\]

Here $m$ denotes a \emph{pointed Kripke model} (a Kripke frame with
valuation plus a designated world), $\varphi$ is a modal formula
(\code{Formula~$\mathbb{N}$} in the Lean source), and $W$ is a multiset
of such pointed models.
The key adequacy result recovers crisp modal truth on singleton
(single-pointed-model) states:

\begin{theorem}[Singleton adequacy]
\label{thm:singleton-kripke}
For any pointed Kripke model $m$ and modal formula $\varphi$,
$m \models \varphi \;\iff\; \Strength(\{m\},\,\varphi) = 1.$
\end{theorem}

\noindent
This is Lean-proven as \code{singleton\_adequacy\_strength\_one} in
\code{PLNWorldModelKripke.lean} (242~lines).

There is also a \emph{consequence lifting} theorem: if
$\forall m.\;m \models \varphi \Rightarrow m \models \psi$ (quantifying
over all pointed Kripke models $m$), then
$\Strength(W,\varphi) \le \Strength(W,\psi)$ for every multiset state
$W$ (a \code{Multiset PointedKripke}).
Singleton semantic implication is equivalent to singleton strength
consequence (\code{pointwiseImplies\_iff\_singletonStrengthLE}).


\subsection{Deontic Kripke consequence rules}
\label{sec:deontic-kripke}

A \code{PointedDeonticKripke} structure extends the above with
labelled-transition-system deontic accessibility. Two modal/deontic
theorems are lifted into WM consequence rules:

\begin{itemize}[leftmargin=2em]
\item \textbf{Reflexive rule:}
  $\Strength(W,\,\mathrm{rexist}(\varphi)) \le \Strength(W,\,\varphi)$.
\item \textbf{DTS obligation--permission:}
  $\Strength(W,\,\mathrm{ob}(\varphi)) \le \Strength(W,\,\mathrm{pe}(\varphi))$.
\end{itemize}

\noindent
Both are packaged as \code{WMConsequenceRule} instances (see
\code{PLNWorldModelKripkeConsequence.lean}, 206~lines). This turns
modal/deontic theorems into \emph{operational WM consequence rules}
rather than leaving them as abstract semantic lemmas.


\subsection{Neighborhood WM instance}
\label{sec:neighborhood-wm}

The neighborhood instantiation is structurally parallel to the Kripke one:
a \code{PointedNeighborhood} pairs a neighborhood model with a world, and
evidence counts neighborhood-satisfaction across a multiset ensemble.
Singleton adequacy and consequence lifting hold by the same schema
(\code{PLNWorldModelNeighborhood.lean}, 274~lines).


\subsection{Neighborhood soundness/completeness bridge}
\label{sec:neighborhood-sc}

The strongest result in this batch connects proof-theoretic derivability
to WM consequence. Given a neighborhood proof system $\mathcal{S}$ and a
frame class $C$ equipped with imported \code{Sound} and \code{Complete}
instances:

\begin{theorem}[Implication-level bridge]
\label{thm:neighborhood-sc}
$\mathcal{S} \vdash (\varphi \to \psi)
\;\iff\;
\mathrm{singletonStrengthLE}_C\,\varphi\,\psi.$
\end{theorem}

\noindent
From provable implication, one obtains a multiset WM strength inequality
over all states whose pointed models lie in the target frame class.

Concrete instantiations are given for $E$, $\mathit{EMN}$,
$\mathit{EMT}$ (reflexive), and $\mathit{ED}$ (serial). See
\code{PLNWorldModelNeighborhoodCompleteness.lean} (315~lines).

An analogous Kripke proof-theoretic bridge is proved in
\code{PLNWorldModelKripkeCompleteness.lean} (238~lines): for a normal
modal logic with Kripke completeness on a frame class $C$,
$\mathcal{S} \vdash (\varphi \to \psi)
\iff \mathrm{singletonStrengthLE}_C\,\varphi\,\psi$
on Kripke WM states. Both semantic regimes now stand at the same
metatheoretic level.

A canonical Kripke-to-neighborhood \emph{representation bridge} is
proved in \code{PLNWorldModelKripkeNeighborhoodCanonical.lean}
(365~lines): each Kripke frame is embedded as a neighborhood frame
via the standard presentation $\mathcal{N}(w)=\{X\mid R[w]\subseteq X\}$,
satisfaction is preserved, and singleton/multiset WM strength agrees
under the embedding. This makes the ``unification'' thesis concrete:
the two instantiations are not merely parallel---they are connected by
a proven transport.

These results are \emph{not} completeness of the generic WM calculus;
they \emph{are} implication-level proof-theoretic closure theorems for
instantiated modal WM semantics. The distinction matters:

\begin{remark}
The generic WM shell remains a \textbf{semantic evaluation framework},
not a standalone logic with a deep state-derivability proof theory.
Every state is trivially derivable
(\code{WMJudgment.axiom}); query judgments are deterministic
(\code{WMQueryJudgment}). The modal instantiations make the shell
\emph{useful} by giving it strong imported semantics.
\end{remark}


\subsection{Governance-to-neighborhood translation}
\label{sec:gov-translation}

The file \code{PLNWorldModelNeighborhoodConsequence.lean} (613~lines)
packages neighborhood consequence rules and provides a recursive
governance-fragment translation:

\begin{center}
\begin{tabular}{ll}
\toprule
\textbf{Governance modality} & \textbf{Neighborhood image} \\
\midrule
$\mathrm{rexist}(\varphi)$ & $\Box\varphi$ \\
$\mathrm{obligatory}(\varphi)$ & $\Box\varphi$ \\
$\mathrm{permitted}(\varphi)$ & $\Diamond\varphi$ \\
$\mathrm{forbidden}(\varphi)$ & $\Box\lnot\varphi$ \\
$\mathrm{optional}(\varphi)$ & $\Diamond\varphi \wedge \Diamond\lnot\varphi$ \\
\bottomrule
\end{tabular}
\end{center}

\noindent
The translation handles propositional connectives recursively and returns
\code{none} for $\mu$/$\nu$ formulas (explicitly outside the current
fragment). All governance action labels on box/diamond operators collapse
to a single neighborhood $\Box$/$\Diamond$. This is a \textbf{disciplined
interoperability layer for a fragment}, not a full faithful embedding of
the governance modal-$\mu$ language.

Modal family theorems are proved for $K$, $C$ (conjunction closure),
McKinsey ($\Box\Diamond\varphi \le \Diamond\Box\varphi$), and Makinson,
with concrete frame class instantiations ($EK$, $EMK$, $EC$, $EMC$).
Endpoint block theorems collect all governance-side inequalities.


\subsection{Philosophical assessment}
\label{sec:modal-philosophy}

Three points emerge from this development.

\paragraph{A.\ The WM calculus is becoming a semantic unifier.}
The WM layer is broad enough to host crisp Kripke semantics on singleton
states, evidential aggregation over multiset ensembles, and neighborhood
semantics with proof-theoretic closure. The right philosophical role for
the WM calculus is not ``this is the final logic of reasoning,'' but
rather: \emph{this is a common evidential/query-calculus interface in
which multiple semantic regimes can live.}

\paragraph{B.\ Multisets of pointed models are an evidential ensemble
semantics.}
A multiset of pointed models reads as a bag of candidate scenarios,
source-indexed hypotheses, or observations with multiplicity. Then
$\Strength(W,\varphi)$ becomes ``how much of the evidential ensemble
supports $\varphi$?'' rather than ``is $\varphi$ true at the actual
world?'' This makes modal semantics relevant to PLN-style evidence
reasoning: it is no longer only about one actual world with
accessibility; it is \emph{structured evidence over candidate worlds}.

\paragraph{C.\ Neighborhood semantics is the more native long-term fit.}
Kripke is the natural first adequacy witness. But neighborhood semantics
supports weaker and non-normal modal behavior (proof systems $E$,
$\mathit{EMN}$, $\mathit{EMT}$, $\mathit{ED}$, etc.), and the
soundness/completeness bridge transports those proof systems into WM
strength inequalities. This is directly relevant if the long-term aim
includes governance, evidence, non-normal modalities, and
context-sensitive support.


\subsection{Limitations}
\label{sec:modal-limits}

For clarity on what is \emph{not} proved:

\begin{enumerate}[leftmargin=2em]
\item The \textbf{generic WM calculus does not have a nontrivial global
  proof theory.} The modal instantiations supply semantics; they do not
  give the shell itself a deep derivability structure.
\item The \textbf{governance translation is partial and lossy}: $\mu$/$\nu$
  formulas map to \code{none}, and all governance action labels collapse.
  It is a bridge for a fragment.
\item The \textbf{base multiset semantics is uniform-count}: each pointed
  model contributes one count.  A \textbf{weighted/source-aware variant}
  is now formalized (\code{PLNWorldModelKripkeWeighted.lean}:
  source-labelled weighted states, trusted-source gating, weight-1
  specialization bridge) and extended to stochastic channels
  (\code{PLNWorldModelExperimentStochastic.lean}: Blackwell utility
  transfer under arbitrary source weights).  The remaining gap is
  \textbf{dependence and correlation correction}: the current weighted
  semantics treats sources as independent.
\end{enumerate}


\subsection{Categorical WM Transport Layer}
\label{sec:categorical-wm-transport}

Beyond modal instantiations, the WM layer now has an explicit categorical
transport surface:
\begin{itemize}[leftmargin=2em]
\item \textbf{Institution transport} over typed WM signatures:
  \code{satisfactionCondition\_evidence},
  \code{satisfactionCondition\_strength},
  \code{satisfactionCondition\_strengthWith},
  \code{satisfactionCondition\_confidence},
  and derived \code{WMInstitution.sat\_condition\_strength}
  (\code{PLNWorldModelInstitution.lean}).
\item \textbf{Hyperdoctrine transport} over typed query fibers with
  change-of-base and Beck-Chevalley:
  \code{beckChevalley\_transport\_exists}
  (\code{PLNWorldModelHyperdoctrine.lean}).
\item \textbf{Unified endpoint}:
  \code{EndpointStatement},
  \code{institution\_beckChevalley\_endpoint},
  \code{endpointSurface\_of\_hyperdoctrine}
  (\code{PLNWorldModelCategoricalBridge.lean}),
  which packages institution-level satisfaction transport and
  hyperdoctrine-level Beck-Chevalley transport in one theorem statement.
\item \textbf{Concrete Ch.8 fixture}:
  \code{ch8\_institution\_beckChevalley\_endpoint\_id\_fixture}
  (\code{PLNWorldModelCategoricalRegression.lean}),
  instantiating the unified endpoint on an identity pullback square.
\item \textbf{Aligned closure wrappers}:
  categorical-surface variants are exposed for the PureKernel$\rightarrow$WM
  bridge and for HOL/FOL WM consequence closures
  (\code{PLNWorldModelPureKernelBridge.lean},
   \code{PLNWorldModelHOLCompleteness.lean},
   \code{PLNWorldModelFOLCompleteness.lean}).
\end{itemize}

For executable regression coverage, the script-level target is
\code{scripts/check\_wm\_categorical.sh}.

\subsection{Stochastic Experiment Layer Theorem Index}
\label{sec:stochastic-experiment-index}

\begin{longtable}{p{0.42\textwidth}p{0.39\textwidth}p{0.15\textwidth}}
\hline
\textbf{Content} & \textbf{Lean theorem(s)} & \textbf{Module} \\ \hline
\endhead
Blackwell factorization utility transport and optimal-value monotonicity &
\code{expectedUtility\_eq\_of\_blackwellFactor},
\code{optimalValue\_mono\_of\_blackwellFactor},
\code{optimalValueWeighted\_mono\_of\_blackwellFactor},
\code{optimalValuePolicy\_mono\_of\_blackwellFactor},
\code{optimalValue\_trustedGate\_mono\_of\_blackwellFactor} &
\code{PLNWorldModelExperimentStochastic} \\
\hline
Concrete finite stochastic fixture endpoints &
\code{fixture\_expectedUtility\_eq\_lifted},
\code{fixture\_optimalValue\_mono},
\code{fixture\_trustedPolicy\_optimalValue\_mono} &
\code{PLNWorldModelExperimentStochasticRegression} \\
\hline
\end{longtable}

Script-level regression target for this layer:
\code{scripts/check\_wm\_experiment.sh}.


%% ============================================================
\section{Related Work}
\label{sec:related}
%% ============================================================

\subsection{Evidence logic and evidence dynamics}

The closest prior work to the WM calculus is the \emph{evidence logic}
program of van Benthem, Fern\'andez-Duque, and
Pacuit~\cite{vanBenthemFernandezDuquePacuit2014,vanBenthemEtAl2013EvidencePlausibility}.
Their \emph{evidence models} $\mathcal{M}=\langle W,E,V\rangle$ assign
to each state a family of evidence sets (neighborhoods). The modal
operator $\Box\varphi$ means ``the agent has evidence for $\varphi$'' and
is non-normal: $\Box\varphi \wedge \Box\psi$ need not imply
$\Box(\varphi\wedge\psi)$, because evidence pieces may be individually
consistent but jointly contradictory. Belief $B\varphi$ is then
\emph{derived} from evidence via a scenario-combination policy. This
architecture---evidence as a first-class structure, queries evaluate
against it, belief/views are derived---is the same pattern as the WM
calculus.

Van Benthem and Pacuit~\cite{vanBenthemPacuit2011} extend this to
\emph{dynamic} evidence management: new evidence is added, old evidence
revised, and belief updates accordingly. Their ``revise evidence, then
extract belief'' pipeline is the closest prior analogue to our
$W \revise \Delta$ followed by $\ev(W,q)$.

The key difference is the evidence type. In evidence logic, evidence
consists of \emph{sets of worlds} (Boolean neighborhoods). In the WM
calculus, evidence is \emph{quantitative}: a pair
$(n^+,n^-)\in\mathbb{R}_{\ge 0}^\infty\times\mathbb{R}_{\ge 0}^\infty$
carrying positive and negative support counts. This makes the WM
calculus oriented toward PLN-style numerical inference rather than
epistemic modal reasoning. But the structural insight---that evidence
should be algebraically rich, not collapsed to a truth value---is shared.

A further distinction: the WM calculus uses \emph{commutative} additive
revision ($W_1 \revise W_2 = W_2 \revise W_1$), reflecting that evidence
aggregation is order-independent. The evidence-dynamics
literature~\cite{vanBenthemPacuit2011} includes richer non-commutative
update operators (public announcement, evidence addition with priority)
where revision order matters. Extending the WM calculus with
order-sensitive revision is future work.

Fitting~\cite{Fitting2017ParaconsistentEvidence} formalizes a
complementary perspective: BLE (Basic Logic of Evidence) allows evidence
to be \emph{incomplete or contradictory}, with a modal embedding into
KX4 where $\Box X$ reads as ``there is implicit evidence for $X$.'' This
structurally parallels the WM's evidence pairs $(n^+,n^-)$ where both
components can be nonzero, though Fitting's framework is qualitative
(justification terms) rather than numeric counts.


\subsection{Neighborhood semantics}

The neighborhood WM instances (\S\ref{sec:neighborhood-wm}--\ref{sec:neighborhood-sc})
build directly on the neighborhood semantics tradition surveyed by
Pacuit~\cite{Pacuit2017NeighborhoodSemantics}. Neighborhood frames
generalize Kripke frames and provide semantics for non-normal modal
logics, including the monotone logics $E$, $\mathit{EMN}$,
$\mathit{EMT}$, and $\mathit{ED}$ for which we prove
soundness/completeness bridges into WM consequence.


\subsection{Provenance semantics and quantale-valued logics}

The WM's evidence quantale and the additivity law
$\ev(W_1 \revise W_2,\, q) = \ev(W_1,q) + \ev(W_2,q)$ (where $+$ is
coordinatewise addition on evidence) place it in the
tradition of \emph{semiring provenance semantics}. Dannert and
Gr\"adel~\cite{DannertGraedel2021SemiringProvenance} show that
semiring-valued model checking for guarded logics yields algebraic
provenance values tracking \emph{how and why} a formula holds---not
merely whether it does. The WM's \code{WMQueryEq} (certified rewrite
producing the same evidence) and \code{evidence\_add} (compositional
aggregation) are instances of this pattern.

Abramsky and Vickers~\cite{AbramskyVickers1993Quantales} establish
quantales as algebraic semantics for logics and processes, unifying
lattice order, monoidal composition, and infinitary join distribution.
This motivates the choice of a quantale (rather than a bare lattice or
semiring) for $\Evidence$: quantale structure is exactly what makes
$\otimes$-distributing-over-joins a theorem rather than an assumption.
Bacci et al.~\cite{BacciMardarePanangadenPlotkin2023} develop
propositional logics interpreted in the Lawvere quantale with proof
systems and decidable completeness results, providing a modern precedent
for quantale-valued views.


%% ============================================================
\section{Roadmap}
\label{sec:roadmap}
%% ============================================================

The most valuable next theorems:

\begin{enumerate}[leftmargin=2em]
\item \textbf{Shell-clarifying theorems.}
Make the meta-status of the generic WM core explicit:
(a)~\emph{triviality}: $\forall W,\; \vdash_{\WM} W$;
(b)~\emph{determinism}: if $\vdash_q W \Downarrow q \mapsto e_1$ and
$\vdash_q W \Downarrow q \mapsto e_2$, then $e_1 = e_2$.
\item \textbf{Weighted / source-labeled multisets.}
Replace plain multisets with source labels, provenance classes, or
weighted multiplicities, so the ensemble semantics models source
reliability and dependence.
\item \textbf{Independence-to-screening-off bridge.}
Prove that d-separation implies the screening-off equalities required by
PLN deduction rules, under explicit positivity.
\item \textbf{VE correctness as a WM-strength rule.}
Show that VE computes the same strength view as querying the BN world
model.
\item \textbf{Quantifier rewrite rules with explicit approximation
  ledgers.}
Hook satisfying-set quantifiers into the rewrite layer with soundness
under stated domain hypotheses (finite, sampled, exchangeable, etc.).
\item \textbf{Operational integration (MORK).}
Implement message passing as inference control that proposes
$\Sigma$-guarded rewrites, while the WM layer remains the single source
of semantic truth.
\end{enumerate}


%% ============================================================
\section*{Acknowledgments}
%% ============================================================

This paper compiles ideas emerging from ongoing collaboration around PLN,
MeTTa, and the Lean formalization in \code{Mettapedia}, with conceptual
inspiration from: Knuth--Skilling foundations of inference, valuation
algebras, belief propagation / junction trees, topos semantics of
quantifiers, and higher-order process calculi.

\bibliographystyle{plain}
\bibliography{references}

\end{document}
